%!TEX program = xelatex
% !TEX encoding = ISO-8859-1
\def\ano{2020}
\def\semestre{1}
\def\disciplina{\'Algebra 1}
\def\turma{C}
\def\numerosemana{02}

\documentclass[12pt]{exam}

\usepackage{caption}
\usepackage{amssymb}
\usepackage{amsmath,amsfonts,amsthm,amstext}
\usepackage[brazil]{babel}
% \usepackage[latin1]{inputenc}
\usepackage{graphicx}
\graphicspath{{/home/jfreitas/GitHub_Repos/Algebra-1/Pictures/}{D:/OneDrive - unb.br/imagens-latex/}}
\usepackage{enumitem}
\usepackage{multicol}
\usepackage{answers}
\usepackage{tikz,ifthen}
\usetikzlibrary{lindenmayersystems}
\usetikzlibrary[shadings]
\Newassociation{solucao}{Solution}{ans}
\newtheorem{exercicio}{}

\setlength{\topmargin}{-1.0in}
\setlength{\oddsidemargin}{0in}
\setlength{\textheight}{10.1in}
\setlength{\textwidth}{6.5in}
\setlength{\baselineskip}{12mm}

\extraheadheight{0.7in}
\firstpageheadrule
\runningheadrule
\lhead{
        \begin{minipage}[c]{1.7cm}
        \includegraphics[width=1.7cm]{unb.pdf}
        \end{minipage}%
        \hspace{0pt}
        \begin{minipage}[c]{4in}
          {Universidade de Bras{\'\i}lia} --
          {Departamento de Matem{\'a}tica}
\end{minipage}
\vspace*{-0.8cm}
}
% \chead{Universidade de Bras{\'\i}lia - Departamento de Matem{\'a}tica}
% \rhead{}
% \vspace*{-2cm}

\extrafootheight{.5in}
\footrule
\lfoot{\disciplina\ - \semestre$^o$/\ano}
\cfoot{}
\rfoot{P\'agina \thepage\ de \numpages}

\newcounter{exercicios}
\renewcommand{\theexercicios}{\arabic{exercicios}}

\newenvironment{questao}[1]{
\refstepcounter{exercicios}
\ifx&#1&
\else
   \label{#1}
\fi
\noindent\textbf{Exerc{\'\i}cio {\theexercicios}:}
}

\newcommand{\resp}[1]{
\noindent{\bf Exerc{\'\i}cio #1: }}
\newcommand{\nsub}{\varsubsetneq}
\newcommand{\vaz}{\emptyset}
\newcommand{\im}{{\rm Im\,}}
\newcommand{\sub}{\subseteq}
\newcommand{\n}{\mathbb{N}}
\newcommand{\z}{\mathbb{Z}}
\newcommand{\rac}{\mathbb{Q}}
\newcommand{\real}{\mathbb{R}}
\newcommand{\complex}{\mathbb{C}}
\newcommand{\cp}[1]{\mathbb{#1}}
\newcommand{\ch}{\mbox{\textrm{car\,}}\nobreak}
\newcommand{\vesp}[1]{\vspace{ #1  cm}}
\newcommand{\compcent}[1]{\vcenter{\hbox{$#1\circ$}}}
\newcommand{\comp}{\mathbin{\mathchoice
{\compcent\scriptstyle}{\compcent\scriptstyle}
{\compcent\scriptscriptstyle}{\compcent\scriptscriptstyle}}}

\begin{document}
    \begin{center}
    {\Large\bf \disciplina\ - Turma \turma\ -- \semestre$^{o}$/\ano} \\ \vspace{9pt} {\large\bf
        Lista de Exerc{\'\i}cios -- Semana \numerosemana}\\ \vspace{9pt} Prof. Jos{\'e} Ant{\^o}nio O. Freitas
    \end{center}
    \hrule

    \vspace{.6cm}

    \begin{center}
        \textit{Observeção: A notação $A \nsub B$ significa que $A \subset B$ e $A \ne B$.}
    \end{center}
    \questao{} Dados os conjuntos $A = \{0,1,2\}$, $B = \{0,2,3\}$ e $C = \{0,1,2,3,4\}$, classifique as afirma\c{c}\~oes a seguir em verdadeira ou falsa, justificando:
    \begin{multicols}{2}
        \begin{enumerate}[label={\alph*})]
            \item $A \sub B$

            \item $A \in C$

            \item $B \sub C$

            \item $\{0,2\} \sub B$

            \item $\{0,2\} \sub C$

            \item $1 \sub C$

            \item $\{1,4\} \in C$

            \item $\{0,3\} \sub B$

            \item $B \sub A$

            \item $\{\vaz\} \sub A$

            \item $\vaz \in C$

            \item $A \sub C$

            \item $x \neq \{x\}$

            \item $\emptyset = \{\emptyset\}$

            \item $\emptyset \in \{\emptyset\}$

            \item $\{1, 3, 3\} = \{1, \{2, 3\}, 3\}$

            \item $\{1, 2, 3, 4, 5\} = \{4, 1, 3, 5, 2, 4, 5\}$
            
            \item $\emptyset \sub \{\emptyset\}$
        \end{enumerate} 
    \end{multicols}

    \questao{} D\^e exemplos de conjuntos n\~ao vazios $A$, $B$ e $C$ tais que:
    \begin{enumerate}[label={\alph*})]
        \item $A \sub B$, $C \sub B$ e $A \cap C = \emptyset$.

        \item $A \sub B$, $C \nsub B$ e $A \cap C = \emptyset$.

        \item $A \sub C$, $A \ne C$ e $B \cap C = \emptyset$.

        \item $A \sub (B \cap C)$, $B \sub C$, $B \ne C$ e $A \ne C$.
    \end{enumerate}

    \vspace{.3cm}

    \questao{} Em cada um dos seguintes itens, determine se a afirma\c{c}\~ao \'e
    verdadeira ou falsa. Se for verdadeira, demonstre-a. Se for falsa, exiba um
    exemplo mostrando que a afirma\c{c}\~ao \'e falsa.
    \begin{enumerate}[label={\alph*})]
        \item Se $x \in A$ e $A \sub B$, ent\~ao $x \in B$.

        \item Se $A \nsub B$ e $B \sub C$, ent\~ao $A \nsub C$.

        \item Se $A \nsub B$ e $B \nsub C$, ent\~ao $A \nsub C$.

        \item Se $x \in A$ e $A \nsub B$, ent\~ao $x \notin B$.

        \item Se $A \sub B$ e $x \notin B$, ent\~ao $x \notin A$.

        \item Se $A$, $B$ e $C$ s\~ao conjuntos, ent\~ao $A \cup (B \cap C) = (A \cup B) \cap C$.
    \end{enumerate}

    \questao{} Demonstre que:
    \begin{enumerate}[label={\alph*})]
        \item Se $A \sub B$ e $C \sub D$, ent\~ao $A \cap C \sub B \cap D$.

        \item $A \cup B = \vaz$ se, e somente se, $A = \vaz$ e $B = \vaz$.

        \item $A \cup B = A \cap B$ se, e somente se,  $A = B$.

        \item $A \sub B$ se, e somente se,  $A \cap B = A$.

        \item $A \cup (B \cap (A \cup C)) = A \cup (B \cap C)$.
    \end{enumerate}

\end{document}